\section{SBL: A Comprehensive Value Accrual Token}

\subsection{The SBL Token Framework}
SBL (Stabolut's native token) is engineered to be more than a simple governance token; it is a comprehensive value accrual vehicle that captures the success of the Stabolut protocol and delivers it to token holders. The SBL framework is built on five key pillars:

\begin{enumerate}
    \item \textbf{Governance Control}: SBL holders have complete authority over the protocol, including treasury allocations, fee structures, risk parameters, and protocol upgrades.
    \item \textbf{Treasury Surplus Distribution}: SBL holders vote on the distribution of excess treasury funds, ensuring that the protocol's success is shared with the community.
    \item \textbf{Progressive Revenue Sharing}: A portion of the protocol's net revenue is distributed to SBL stakers, creating a direct link between protocol usage and token holder returns.
    \item \textbf{Deflationary Mechanisms}: The protocol can use treasury surplus to buy back and burn SBL tokens, reducing the circulating supply and increasing the value of the remaining tokens.
    \item \textbf{Utility Benefits}: SBL holders receive tangible benefits within the Stabolut ecosystem, such as reduced fees and enhanced voting power.
\end{enumerate}

\subsection{Progressive Staking and Governance}
To incentivize long-term commitment and active participation, Stabolut implements a progressive staking system with three tiers:

\begin{itemize}
    \item \textbf{Tier 1 (30-day lock)}: 1.2x voting multiplier, 5\% fee discount.
    \item \textbf{Tier 2 (180-day lock)}: 1.5x voting multiplier, 15\% fee discount.
    \item \textbf{Tier 3 (365-day lock)}: 2.0x voting multiplier, 30\% fee discount.
\end{itemize}

This tiered system ensures that the most committed community members have the greatest influence on the protocol's direction and receive the most significant benefits.

\subsection{Treasury Surplus Distribution: A Democratic Approach to Value Creation}
The cornerstone of the SBL value accrual framework is the governance-controlled treasury surplus distribution mechanism. This mechanism is activated when the treasury's reserves exceed 50\% of the circulating USB supply plus a buffer for operational expenses.

When this threshold is met, SBL holders can vote on how to allocate the surplus funds. The available options are:

\begin{itemize}
    \item \textbf{Direct Airdrop}: A proportional distribution of the surplus funds to SBL holders.
    \item \textbf{Buyback and Burn}: The protocol uses the surplus to purchase SBL tokens on the open market and permanently removes them from circulation.
    \item \textbf{Yield Enhancement}: The surplus is deployed into additional yield-generating strategies, with the returns passed on to SBL stakers.
\end{itemize}

This democratic approach to value distribution ensures that the community has direct control over the protocol's economic engine, aligning the interests of the protocol with those of its token holders.

\subsection{Regulatory Compliance}
The SBL value accrual framework is designed to be fully compliant with evolving regulatory landscapes. By structuring SBL as a utility token with governance rights and making distributions non-guaranteed and community-controlled, Stabolut mitigates the risk of SBL being classified as a security. All major governance decisions are subject to legal review to ensure ongoing compliance.

